\section{Revealed-Sacrifice Observation Model}
\label{sec:sacrifice_model}

The privacy-compatible observation channel for $\volR$ is
\emph{revealed sacrifice}: agents disclose the value they place on
unbenchmarked capabilities through the benchmarked capabilities
they voluntarily surrender to obtain them.  The framework does not
observe \emph{what} the agent values privately---only \emph{what
they commit to publicly through sacrifice}.

\subsection{Core definitions}

\begin{definition}[Revealed-Sacrifice Event]
\label{def:sacrifice_event}
A revealed-sacrifice event is a tuple $(i, X, Y, t)$ where:
\begin{itemize}
\item $i$ is the agent performing the sacrifice;
\item $X$ is a benchmarked capability (or capability bundle) with
  known $\Delta\volL(X) \geq 0$, surrendered by agent~$i$ at
  time~$t$;
\item $Y = \{Y_1, \ldots, Y_k\}$ is an unbenchmarked capability
  bundle acquired by agent~$i$;
\item $t$ is the event timestamp.
\end{itemize}
The event emits the signal:
\[
  \volRlower(Y) \;\geq\; \Delta\volL(X)
\]
under the assumptions S0 (calibration), S1 (free choice), and S2
(benchmark grounding) stated below.
\end{definition}

\paragraph{Justification (revealed-preference inequality).}
The transfer from the microeconomic revealed-preference
inequality~\citep{samuelson1938note}
$\Ui(Y) \geq \Ui(X)$ to the framework-level bound
$\volRlower(Y) \geq \Delta\volL(X)$ requires two bridge
conditions:

\begin{enumerate}
\item \textbf{Sacrifice-side grounding (Assumption S2):} The
  agent's valuation of the benchmarked sacrifice $X$ is calibrated
  to its $\volL$ contribution:
  $\Ui(X) \geq \Delta\volL(X)$.  This holds by construction when
  $\volL$ is the operational target on the benchmarked subspace.

\item \textbf{Acquisition-side calibration (Assumption S0):} The
  agent's valuation of the unbenchmarked bundle $Y$ is a lower
  bound on $Y$'s $\volR$ contribution:
  $\Ui(Y) \leq \volR(Y)$.  This asserts that agents do not
  systematically overvalue unbenchmarked bundles relative to their
  exercise contribution.
\end{enumerate}

Given both conditions:
$\volR(Y) \geq \Ui(Y) \geq \Ui(X) \geq \Delta\volL(X)$,
completing the chain.

\begin{remark}[S0 is the weakest calibration assumption]
S0 does \emph{not} require that agent valuations perfectly track
$\volR$---only that they do not \emph{overstate} it.  Agents may
undervalue unbenchmarked bundles (sacrifice little for something
very valuable), in which case the lower bound is conservative but
still valid.  S0 fails only in the overpayment direction.
\end{remark}

\begin{remark}[S1 filters below-sufficiency sacrifices]
\label{rem:s1_filters}
The need-sufficiency architecture
(Section~\ref{sec:need_sufficiency},
Proposition~\ref{prop:s1_sufficiency}) establishes that
below-sufficiency sacrifices---events where the agent trades to
meet a need on which they are below the sufficiency
threshold---fail S1's non-degrading alternative condition.  These
events do not enter the sacrifice aggregate.  The sacrifice channel
therefore operates only over events where S1 holds, and the
framework accounts for below-sufficiency sacrifice evidence
separately via the need-sufficiency diagnostic
(Section~\ref{sec:gap_decomposition}).
\end{remark}

\begin{definition}[Aggregate Trade Window]
\label{def:trade_window}
An aggregate trade window $[t_0, t_0 + T]$ is an observation
period over which revealed-sacrifice events are collected.  The
window parameters are:
\begin{itemize}
\item $T$: window duration;
\item $N$: number of observed sacrifice events in the window;
\item $\{(i_n, X_n, Y_n, t_n)\}_{n=1}^N$: the event sequence.
\end{itemize}
The aggregate trade window is the temporal analogue of the
observation window $\Delta$
in~\citep[Definition~4]{lasser2026scm}---it defines the timescale
over which the framework accumulates sacrifice evidence.
\end{definition}

\subsection{Privacy properties}

\begin{theorem}[Privacy-Minimality of Revealed-Sacrifice
Observation]
\label{thm:privacy}
The revealed-sacrifice channel satisfies the following five
privacy properties, each of which is load-bearing for
compatibility with the framework's prior commitments:

\begin{enumerate}
\item[\textbf{(P1)}] \textbf{No interior access.}  The framework
  never observes what the agent values privately---only what they
  commit to publicly through sacrifice.  The observation channel
  records the tuple $(i, X, Y, t)$, not the agent's internal
  valuation function $\Ui$.

\item[\textbf{(P2)}] \textbf{Consent via participation.}  Markets
  and labour exchanges are opt-in.  Non-trading is honoured by
  silence: an agent who does not participate in any trade emits no
  signal, and the framework draws no inference about that agent's
  $\volR$ contribution.

\item[\textbf{(P3)}] \textbf{Discretization as a feature.}  Trade
  events are sparse by construction.  The signal bandwidth is
  bounded by trade frequency, not by observation intensity.

\item[\textbf{(P4)}] \textbf{Commitment-layer composability.}
  The commitment protocol
  of~\citep[Definition~6]{lasser2026exo} extends to
  revealed-sacrifice events: a sacrifice event can be committed as
  a proof-of-trade exposing only the $\volR$-category and
  $\Delta\volL$-magnitude, not counterparty, price, or specific
  goods (Section~\ref{sec:commitment}).

\item[\textbf{(P5)}] \textbf{Third-party observability (money
  channel and public time channel).}  Unlike controlled
  relaxation~\citep{lasser2026wamura} (which requires the
  framework to \emph{impose} a test) and unlike
  verification~\citep{lasser2026exo} (which requires the agent to
  \emph{participate} in a commitment), money-sacrifice events and
  publicly observable time-sacrifice events are observable from
  public-facing ledger data (market exchanges, payment rails,
  public commitments) with no agent participation required.  P5
  does not extend to private time-sacrifice activities (parenting,
  contemplation), which require voluntary disclosure.
\end{enumerate}
\end{theorem}

\begin{proof}[Proof sketch]
(P1):~By construction---the observation channel records
$(i, X, Y, t)$, not $\Ui$.
(P2):~The lower bound sums over observed events; an agent with
zero events contributes zero.
(P3):~Channel bandwidth is $|\mathrm{events}|/T$, bounded by
trade frequency.
(P4):~Composition with~\citep[Definition~6]{lasser2026exo}; the
commitment binds $(\Delta\volL(X), \catY)$ without revealing
specifics (see Section~\ref{sec:commitment}).
(P5):~Market transactions leave observable traces even without
agent cooperation; the adversarial model is strictly weaker than
that of~\citep{lasser2026wamura}.
\end{proof}
